\section{Pengantar Kuantifikasi dan Semesta Pembicaraan}

Sebagaimana dibahas pada Bab 6, fungsi proposisional (misalnya $P(x) \equiv \text{``}x > 3\text{''}$) belum memiliki nilai kebenaran tetap karena masih memuat peubah bebas. Agar menjadi proposisi (kalimat tertutup yang bernilai benar atau salah), peubah dapat disubstitusi dengan konstanta tertentu. Selain substitusi langsung, tersedia pendekatan yang lebih umum dan ekspresif, yaitu \logic{kuantifikasi} \cite{rosen2019}.

\subsection{Kuantifikasi}

\begin{definisi}[Kuantifikasi]
\logic{Kuantifikasi} (\textit{quantification}) adalah proses pemberian kuantor pada fungsi proposisional untuk menyatakan jangkauan objek yang dibicarakan. Operator yang digunakan disebut \logic{kuantor} (\textit{quantifier}).
\end{definisi}

Melalui kuantifikasi, pernyataan tidak lagi terbatas pada satu objek tertentu, tetapi dapat menyatakan sifat untuk seluruh elemen domain atau untuk sebagian elemen domain. Kerangka formal yang mempelajari fungsi proposisional beserta kuantornya adalah \logic{Logika Predikat} atau \logic{Logika Orde Pertama} (\textit{First-Order Logic}).

\subsection{Semesta Pembicaraan (Universe of Discourse)}

Sebelum mengevaluasi kalimat berkuantor, semesta pembicaraan harus ditentukan secara eksplisit.

\begin{definisi}[Semesta Pembicaraan]
\logic{Semesta Pembicaraan} (\textit{universe of discourse} atau \textit{domain of discourse}), disimbolkan dengan $\mathcal{U}$, adalah himpunan semua nilai yang diperbolehkan untuk menggantikan variabel pada suatu fungsi proposisional.
\end{definisi}

Pemilihan $\mathcal{U}$ bersifat menentukan, karena nilai kebenaran proposisi berkuantor dapat berubah ketika domain berubah.

\begin{contoh}
\textbf{Pengaruh semesta pembicaraan}

Diberikan pernyataan: ``Semua elemen $x$ memenuhi $x^2 > 0$''.
\begin{itemize}
    \item Jika $\mathcal{U} = \mathbb{R}$, pernyataan bernilai \textbf{Salah} karena $x=0$ menghasilkan $0^2=0$, dan $0 \not> 0$.
    \item Jika $\mathcal{U} = \mathbb{Z}^+$, pernyataan bernilai \textbf{Benar} karena kuadrat bilangan bulat positif selalu positif.
\end{itemize}
\end{contoh}

\subsection{Kategori Kuantor Utama}

Dalam logika predikat terdapat dua kuantor utama:
\begin{enumerate}
    \item \textbf{Kuantor universal} (\textit{universal quantifier}), disimbolkan dengan ``$\forall$'', untuk menyatakan bahwa suatu properti berlaku bagi seluruh elemen domain.
    \item \textbf{Kuantor eksistensial} (\textit{existential quantifier}), disimbolkan dengan ``$\exists$'', untuk menyatakan bahwa ada sekurang-kurangnya satu elemen domain yang memenuhi properti tertentu.
\end{enumerate}

Pada sebagian literatur \cite{wiki_quantifier_logic}, domain ditulis eksplisit di bawah kuantor, misalnya $\forall_{x \in \mathbb{Z}} P(x)$. Di buku ini, apabila domain tidak ditulis eksplisit, domain diasumsikan telah ditetapkan oleh konteks.
